\subsection{Kleine Gremienkunde}

Des Deutschen liebstes Kind ist -- nein, nicht sein Auto! Die B"urokratie, denn 
ohne sie herrschte Chaos im Dunkel und Angst, Furcht und Schrecken all"uberall. 
Damit auch die Studierenden sich gut verwaltet f"uhlen d"urfen, gibt es 
nat"urlich ebenfalls an der TU Braunschweig eine Menge Gremien und Organe, die 
Entscheidungen f"allen (oder verschieben ;-)), Kompetenzen zuteilen (oder 
verschieben ;-)) und aufpassen, dass alles mehr oder weniger seinen 
demokratischen Gang geht.

\begin{figure}[h]
  \centering\includegraphics[width = 0.7\linewidth]{bilder/comics/otto1_1.png}
\end{figure}

Damit ihr euch im Dschungel ein wenig besser orientieren k"onnt, wollen wir im 
Folgenden versuchen, die einzelnen Gremien und deren Aufgaben vorzustellen und 
euch zeigen, wie und in welchem Umfang ihr unmittelbar (durch Wahl) oder 
mittelbar (durch die Gew"ahlten) Einfluss auf die Hochschulpolitik nehmen 
k"onnt. F"ur die, die lieber Bildchen begucken, stellen wir eine Grafik zur 
Verf"ugung, die ihr auf der Seite vor diesem Artikel findet -- schenkt ihr 
Beachtung, sie hat es verdient!

Lange Rede, gar kein Sinn, wir fangen an:

\subsubsection*{Organe der Studierendenschaft}

Dieser inzwischen von allen maskulinen und femininen Kennzeichen befreite 
Begriff vereint nichts anderes als alle StudentInnen der TU-BS unter sich (also 
auch DICH!). Die StudentInnen, die mehr oder weniger zuf"allig an der gleichen 
Fakult"at studieren, fasst man als \textbf{Fachschaft (FS)} zusammen, derer 
gibt es 10 an der guten alten Carolo-Wilhelmina. Eigentlich solltest du es 
inzwischen mitbekommen haben, aber du, verehrte/r LeserIn, geh"orst mit gro"ser 
Wahrscheinlichkeit zur Fachschaft "`Mathematik und Informatik"'. Gibt es 
innerhalb einer Fachschaft noch Unterschiede in den Studienrichtungen, so wird 
in \textbf{Fachgruppen (FG)} aufgeteilt, bei uns ist das die Fachgruppe 
Mathematik und die Fachgruppe Informatik.

Die Studierenden einer Fachschaft werden "ublicherweise durch den 
\textbf{Fachschaftsrat (FSR)} vertreten. Hier hat sich schon seit Jahren eine
Extrawurst-Regelung f"ur unsere Fakult"at eingeb"urgert: Da die 
Studienf"acher Mathematik und Informatik doch sehr unterschiedlich sind und die
notwendigen Absprachen immer ausreichend informell getroffen werden konnten, 
gibt es keinen Fachschaftsrat "`Mathematik und Informatik"'. Stattdessen
konzentriert sich die Energie auf die \textbf{Fachgruppenr"ate (FGR)}, in eurem 
Fall den Fachgruppenrat Informatik. Er k"ummert sich um die Belange der
Fachgruppe, beruft die Fachgruppen-Vollversammlungen ein, streitet sich mit der
Fakult"at, wenn's mal wieder Meinungsverschiedenheiten wegen irgendwelcher 
Neuerungen gibt, organisiert die Orientierungseinheit f"ur die Erstsemester am 
Anfang des Wintersemesters, verleiht Pr"ufungsprotokolle, informiert durch 
Herausgeben der Fachgruppenzeitschrift \nte\ und "uber das Internet 
\nurlfootnote{http://fginfo.cs.tu-bs.de/} und tr"agt das ganze Semester "uber 
Informationen aus den verschiedenen Gremien zusammen, und an euch weiter.
F"ur dich ist der FGR der 
wichtigste Ansprechpartner, denn auch wenn wir deine Probleme mal nicht l"osen 
k"onnen, dann k"onnen wir dir wenigstens sagen, an wen oder was du dich wenden 
kannst. Damit auch zwischen den verschiedenen Fachschaften
und Fachgruppen kommuniziert wird, 
gibt es das \textbf{Fachschaftenplenum}, was kein Gremium im eigentlichen Sinne 
ist, aber ein Forum zum Meinungs- und Interessenaustausch darstellt. Es trifft 
sich etwa einmal im Monat und ist f"ur jeden offen, der einen Einstieg in die 
Unipolitik sucht.

Ganz basisdemokratisch ist auf allen Hierarchie\-ebenen der Studierendenschaft
die jeweilige \textbf{Vollversammlung (VV)} das oberste Organ, allerdings nur
mit empfehlendem Character. Sie findet ein- bis zweimal pro Jahr statt und
dort wird "uber Aktuelles und Wichtiges informiert und/oder abgestimmt. Eine
Vollversammlung aller Studierenden wird vom StuPa-Pr"asidium, eine 
Fachschafts- oder Fachgruppen-VV vom FSR oder FGR einberufen und geleitet.

Womit wir bei Abk"urzungen w"aren, die noch nicht erkl"art wurden -- aber keine 
Bange, das kommt: Das \textbf{Studierendenparlament (StuPa, SP)} ist die 
unmittelbare Vertretung aller StudentInnen und wird von der Studierendenschaft 
direkt in jedem Semester gew"ahlt, und tagt \textbf{hochschul"offentlich}.
Die etwa 40 Mitglieder des StuPa 
beschlie"sen studentische Angelegenheiten, verabschieden den studentischen 
Haushalt und w"ahlen den \textbf{Allgemeinen Studierendenausschuss (AStA)},
den \textbf{"Ubergeordneten Wahlausschuss ("UgWA)}
und verschiedene weitere Aussch"usse. Das StuPa w"ahlt au"serdem sein eigenes
Pr"asidium, welches die Sitzungen und (uniweiten) Vollversammlungen leitet und
das StuPa nach au"sen hin
vertritt.

\begin{figure}[h]
  \centering\includegraphics[width = 0.7\linewidth]{bilder/comics/otto1_2.png}
\end{figure}

Von allen studentischen Aussch"ussen ist der \textbf{AStA} sicherlich der
sichtbarste. Er ist das ausf"uhrende Organ der 
Studierendenschaft und vertritt alle Studierenden nach au"sen, z.B. bei 
Verhandlungen mit der BVAG wegen des Semestertickets. Seine Aufgaben werden vom 
StuPa festgelegt und beinhalten z.B. Serviceangebote (Kopieren, Binden, 
Internationaler Studiausweis) oder Informationsquellen zu den 
unterschiedlichsten Themen. Um sich zu entlasten, kann er ReferentInnen 
bestellen, die sich um einzelne Bereiche mehr oder weniger hauptamtlich 
k"ummern. Das zweite vom StuPa gew"ahlte Gremium ist der \textbf{"Ubergeordnete 
Wahlausschuss ("UgWA)}, der die studentischen Wahlen organisiert und "uberwacht.

\subsubsection*{Kollegialorgane}

Neben den bis jetzt vorgestellten Organen der Verfassten Studierendenschaft 
gibt es nat"urlich auch noch Schnittstellen zwischen den Studis und den anderen 
an der Universit"at vertretenen Personengruppen, den MTVlern (MitarbeiterInnen 
aus Technik und Verwaltung), den WiMis (Wissenschaftliche MitarbeiterInnen, 
AssistentInnen) und nat"urlich den Lehrenden (ProfessorInnen). Hier ist das 
oberste Organ innerhalb der Fakult"aten der \textbf{Fakult"atsrat (FKR)}, 
dem 7 Professoren, 2 Studis, 2 MTVler und 2 WiMis angeh"oren. Hier wird all das 
entschieden, was andere Gremien oder das Dekanat erarbeitet haben, bspw. 
"Anderungen an der BPO. Wird eine Entscheidung getroffen, so ist diese 
sozusagen offiziell geworden und kann umgesetzt werden. Da auf Grund der 
Stimmenverteilung (s.o.) die Professoren immer eine Mehrheit haben, m"ussen wir 
in den Gremien, die vorher die inhaltliche Arbeit leisten, versuchen, unsere 
und eure Vorstellungen einzubringen. Die studentischen 
Vertreter werden einmal im Jahr, jeweils im Wintersemester, direkt gew"ahlt. Da 
wie gesagt die Mathematik und die Informatik doch durchaus unterschiedliche 
Studieng"ange sind, gibt es einen nicht formellen "`kleinen Fakult"atsrat"', 
die \textbf{Informatik-Kommission}. Die Informatik-Kommission, die im 
Verh"altnis 3 : 1 : 1 : 1 besetzt ist, ber"at informatikspezifische Dinge und 
bereitet sie f"ur den Fakult"atsrat vor, damit die Entscheidungen im FKR 
schneller gef"allt werden k"onnen und sich die Mathematiker nicht so langweilen 
;-).

Das formal oberste Gremium der Uni ist der \textbf{Senat}, der sich mit 
allgemeinen Sachen befasst, die "uber der Zust"andigkeit der Fakult"aten 
liegen (als wichtiger Punkt ist hier die Verteilung des universit"aren 
Haushaltes zu nennen). Wie in den FKR ist hier die Stimmengewichtung 7 : 2 : 2 
: 2, auch seine Mitglieder werden j"ahrlich gew"ahlt. Wie der AStA hat auch der 
Senat die M"oglichkeit, seine Arbeit unterst"utzende Kommissionen einzusetzen.

\subsubsection*{Kommissionen und Aussch"usse}

Da wir so oft Kommissionen und Aussch"usse erw"ahnt haben, seien die drei 
wichtigsten hier kurz vorgestellt: zun"achst ist da die 
\textbf{Studienkommission (StuKo)} zu erw"ahnen, die mit dem neuen 
Nieders"achsischen Hochschulgesetz (NHG) im vergangenen Jahr eingef"uhrt wurde. 
Sie ist das einzige gemischte Gremium, in dem die Studierenden die Mehrheit 
haben: 1 : 2 : 0 : 1 lautet die Verteilung der stimmberechtigen Mitglieder. Die 
Studienkommission erarbeitet vor allem Vorschl"age f"ur die Verbesserung der 
Qualit"at in der Lehre, so werden z.B. Vorschl"age zur "Anderung der 
Studienordnung und der BPO diskutiert. Die Studienkommission muss vor allen 
Entscheidungen des Fakult"atsrates, welche die Lehre, das Studium oder 
Pr"ufungen betreffen, angeh"ort werden. Eingesetzt wird die StuKo von den 
Fakult"atsr"aten, die studentischen Vertreter rekrutieren sich meist aus den 
FSR/FGRn oder deren Umfeld (obwohl theoretisch jede/r Interessierte mitarbeiten 
kann). Die Sitzungen sind hochschul"offentlich, d.h. auch nicht gew"ahlte 
Studierende k"onnen (und sollten) dort jederzeit ihre Stimme einbringen.

\begin{figure}[h]
  \centering\includegraphics[width = 0.7\linewidth]{bilder/comics/otto1_3.png}
\end{figure}

Auch Professoren ist es einmal verg"onnt, sich in den Ruhestand zu begeben oder 
andere Hochschulluft zu schnuppern. Wenn dies ansteht, dann muss die 
freigewordene Stelle (logischerweise) in den meisten F"allen neu besetzt 
werden. Daf"ur wird eine \textbf{Berufungskommission} vom Senat eingesetzt, um 
die Nachfolge zu regeln. Hier werden die Kandidaten, nachdem eine Vorauswahl 
getroffen wurde, sozusagen auf Herz und Nieren "uberpr"uft, und zwar im Rahmen 
eines "offentlichen Vortrags, den sich jede/r Interessierte anh"oren kann. Die 
zwei studentischen Vertreter in der Kommission interessiert dabei vor allem, ob 
der/die KandidatIn f"ahig ist, eine Vorlesung verst"andlich und klar 
strukturiert zu halten oder ob er sich in schweren wissenschaftlichen 
Formulierungen verliert, denn es gibt immer wieder Personen, die sich
haupts"achlich auf die Forschungs- und kaum auf die Lehraufgaben konzentrieren.
Die Berufungskommission 
erstellt nach ausgiebigen Beratungen eine Liste, die, nachdem sie den Senat 
passiert hat, ans Kultusministerium (Verzeihung, "`Ministerium f"ur 
Wissenschaft und Kultur"' (MWK)) weitergeleitet wird, das dann nach dieser 
Liste entscheidet, mit wem es, vertreten durch den Uni-Pr"asidenten, der ja 
formal auch Angestellter des MWK ist, in Verhandlungen tritt.

Ein ziemlich wichtiger, von den FKR eingesetzter Ausschuss ist der 
\textbf{Pr"ufungsausschuss (PA)}. Der PA besteht aus 5 Mitgliedern (3 : 1 : 0 : 
1)
und ist f"ur alle Fragen zust"andig, die im Zusammenhang mit Pr"ufungen
auftreten k"onnen. Bei (fast) allen Problemen, die mit Pr"ufungen 
zusammenh"angen, kann
man sich an den Pr"ufungsausschuss wenden - so k"onnen z.B. weitere
Nebenf"acher auf Antrag der Studierenden vom Pr"ufungsausschuss genehmigt
werden. 

Und -- last but not least -- sei die \textbf{Frauenversammlung} erw"ahnt, das 
einzige Gremium mit Stimmengleichheit (10 : 10 : 10 : 10). Sie wird von allen 
weiblichen Studentinnen und Mitarbeiterinnen gew"ahlt und bestimmt aus ihren 
Reihen die universit"are Frauenbeauftragte, die sich f"ur Gleichstellung und 
-berechtigung der Frauen an der Uni einsetzt. Sie "uberwacht beispielsweise, ob 
in den einzelnen Aussch"ussen auch Frauen vertreten sind, ob Frauen in 
irgendeiner Art und Weise diskriminiert werden oder ob die gesetzlichen 
Frauenquoten in den "Amtern eingehalten werden. 

Daneben gibt es nat"urlich noch ungez"ahlte weitere kleine und gro"se Gremien, 
Aussch"usse, Kommissionen und damit verbunden viele viele P"ostchen, die immer 
wieder zu vergeben sind. Wenn ihr also Blut geleckt habt und nicht nur durch 
eure Beteiligung bei den Wahlen Einfluss auf die Hochschulpolitik nehmen wollt, 
dann meldet euch doch im Fachschaftsrat und arbeitet mit -- ihr seid herzlich 
willkommen!

\nauthor{\copyright{} Fachschaftsrat Maschinenbau, FGR Informatik\\WS, PK}
